\chapter{Brownian Motion}

In Chapter~28, we saw how martingales balance fair games at discrete clock
ticks, keeping future expectations anchored to where you stand today. What happens
if you run that fair game continuously, flipping the coin a million times a second
while shrinking each step to microscopic size?

In the summer of 1827 the Scottish botanist Robert Brown put a drop of water
under his microscope. Suspended in the water were tiny grains of pollen, and
inside the grains were even tinier particles. Brown watched these particles and
saw something unexpected: they would not sit still. They jittered, darted,
and moved irregularly in ceaseless motion. His first thought was that they
were alive. But he tested grains from a plant pressed in a herbarium a century
earlier, then specks of ground glass, then mineral dust -- and
they all moved in the same erratic way. Whatever was making them move, it was not life. It was
a property of the liquid itself.

It took most of a century to understand what Brown had seen. The answer, when it
came, tied together three key ideas in science: that matter is made of
molecules, that heat is molecular motion, and that the visible jitter of a pollen
grain is the accumulated effect of billions of invisible water molecules colliding
with it from every side. The particle undergoes an uninterrupted barrage of collisions,
producing a random path.

This chapter is about the \emph{mathematical} object that models this erratic motion. We
will build it the way nature builds it: as the limit of a random walk taken with
smaller and smaller, faster and faster steps. The limit is
\emph{Brownian motion}, also called the \emph{Wiener process}.
It is continuous everywhere yet differentiable nowhere.
Its increments are normal. It carries the Central Limit Theorem over into continuous-time
randomness. It solves the heat equation, and it underpins
the Black--Scholes formula in mathematical finance. In this chapter, we develop both the
analytical properties and the geometric intuition for Brownian motion.

\chapterline

\section{From Random Walk to a Wild Curve}
\label{ch29:walk-to-curve}

You have already met the simple random walk. A walker stands at the origin. Every
unit of time she flips a fair coin and steps $+1$ or $-1$. After $n$ steps her
position is
\[
  S_n = X_1 + X_2 + \dots + X_n,
  \qquad X_i = \pm 1 \text{ each with probability } \tfrac12.
\]
The steps are independent, each has mean $0$ and variance $1$, so
\[
  \E[S_n] = 0, \qquad \Var(S_n) = n.
\]
The walk drifts nowhere on average, but it spreads: its typical distance from the
origin grows like $\sqrt{n}$, the square root of the number of steps. This
$\sqrt{n}$ scaling governs the entire chapter.

\begin{intuitionbox}
Why $\sqrt n$ and not $n$? Because the steps point in random directions and
mostly cancel. If they all pointed the same way you would travel a distance $n$.
If they were perfectly anti-correlated you would never leave home. Independent
random signs land you in between, at the scale $\sqrt n$. The variance adds (that
is the one thing variance always does for independent terms), and standard
deviation is its square root. Diffusion is slow precisely because it is built out
of cancellation.
\end{intuitionbox}

Now we want a \emph{continuous} version. The trick is a clever rescaling. Suppose
we want to watch the walk over the time interval $[0,1]$, but take $n$ steps in
that interval instead of one. Then each step lasts time $1/n$. To keep the spread
finite as $n\to\infty$ we must also shrink each spatial step. How much? Variance
must stay finite. After $n$ steps the variance is $n$ times the variance of one
step. If each step has size $\Delta x$, its variance is $(\Delta x)^2$, so the
total variance is $n(\Delta x)^2$. For this to equal $1$ (one unit of time, one
unit of variance) we need
\[
  (\Delta x)^2 = \frac1n,
  \qquad\text{i.e.}\qquad \Delta x = \frac{1}{\sqrt n}.
\]
Here is the scaling that makes everything work: \textbf{time shrinks like
$1/n$, but space shrinks only like $1/\sqrt n$.} Space and time scale
\emph{differently}. This asymmetry -- space as the square root of time -- is the
signature of diffusion, and it is why Brownian motion looks the way it does.

Define the rescaled walk
\[
  W^{(n)}(t) \;=\; \frac{1}{\sqrt n}\, S_{\floor{nt}},
  \qquad 0 \le t \le 1,
\]
interpolated linearly between the integer step times. At $t=1$ this is
$S_n/\sqrt n$, which by the Central Limit Theorem converges to a standard normal.
But the CLT only tells you about the endpoint. The stronger statement is that the
\emph{entire path} $t\mapsto W^{(n)}(t)$ converges, as a random function, to a
limiting random function. That limit is Brownian motion. Mathematicians call this
path-level Central Limit Theorem \emph{Donsker's invariance principle}: whether
the microscopic steps come from coin flips, dice rolls, or continuous noise, as
long as they are independent with mean zero and finite variance, their trajectories
converge to the same continuous limit.

\begin{figure}[H]
\centering
\begin{tikzpicture}
\begin{axis}[
  width=12cm, height=6.2cm,
  xlabel={time $t$}, ylabel={position},
  title style={font=\small},
  title={Rescaled random walks approach a continuous path},
  legend pos=north west, legend cell align=left,
  grid=both, grid style={gray!15},
  ymin=-0.95, ymax=1.45,
]
\addplot[sourceblue, thick] coordinates {
 (0,0)(0.0625,0.25)(0.125,0)(0.1875,0.25)(0.25,0.5)(0.3125,0.25)
 (0.375,0.5)(0.4375,0.25)(0.5,0)(0.5625,-0.25)(0.625,0)(0.6875,-0.25)
 (0.75,-0.5)(0.8125,-0.25)(0.875,-0.5)(0.9375,-0.75)(1,-0.5)};
\addlegendentry{$n=16$ (coarse)}
\addplot[quotientorange, thick] coordinates {
 (0,0)(0.05,0.18)(0.1,0.05)(0.15,0.30)(0.2,0.42)(0.25,0.28)(0.3,0.50)
 (0.35,0.38)(0.4,0.55)(0.45,0.40)(0.5,0.22)(0.55,0.05)(0.6,0.20)
 (0.65,-0.02)(0.7,-0.25)(0.75,-0.40)(0.8,-0.22)(0.85,-0.45)(0.9,-0.60)
 (0.95,-0.48)(1,-0.55)};
\addlegendentry{$n=64$ (finer)}
\addplot[densityred, very thick] coordinates {
 (0,0)(0.04,0.14)(0.08,0.09)(0.12,0.26)(0.16,0.20)(0.2,0.38)(0.24,0.30)
 (0.28,0.44)(0.32,0.40)(0.36,0.52)(0.4,0.47)(0.44,0.36)(0.48,0.27)
 (0.52,0.15)(0.56,0.10)(0.6,0.18)(0.64,0.05)(0.68,-0.10)(0.72,-0.22)
 (0.76,-0.34)(0.8,-0.28)(0.84,-0.40)(0.88,-0.52)(0.92,-0.55)(0.96,-0.50)(1,-0.55)};
\addlegendentry{large $n$ (Brownian)}
\end{axis}
\end{tikzpicture}
\caption{As we take more, smaller steps with the diffusive scaling
$\Delta x = 1/\sqrt n$, the jagged random walk fills in to a continuous but
infinitely wiggly curve. The endpoint distribution stays fixed (standard normal);
what converges is the whole path.}
\label{ch29:fig-scaling}
\end{figure}

\begin{historybox}
Brown observed the jitter in 1827 and published his account in 1828, showing
the motion was non-biological, though he had no physical explanation. In 1900,
Louis Bachelier modeled stock prices with the same random-walk limit in his doctoral
thesis.

The physical explanation arrived in 1905: Albert Einstein
derived the law $\E[\text{displacement}^2] = 2Dt$ from the hypothesis that heat is
molecular motion. The diffusion constant $D$ depended directly on Avogadro's number,
offering the first quantitative test for the physical reality of atoms. Marian
Smoluchowski independently derived the same law in 1906. In 1908, Jean Perrin
measured the jitter under a microscope, confirmed Einstein's formula, and won the
1926 Nobel Prize in Physics for putting the existence of molecules beyond doubt.

The rigorous mathematical foundation -- constructing a probability measure on the
space of continuous, nowhere-differentiable paths -- was completed by Norbert Wiener
in 1923, which is why Brownian motion is called the \emph{Wiener process} and denoted
$W(t)$. In the 1940s, Kiyosi It\^o developed stochastic calculus to integrate along
these paths.
\end{historybox}

\section{The Defining Properties}
\label{ch29:definition}

Rather than chase the limit through all its technical thickets, we do what
mathematicians eventually learned to do: we \emph{list the properties} the limit
must have and take \emph{those} as the definition. Wiener proved that a process
with these properties exists and is essentially unique.

\begin{definition}[Standard Brownian motion]
\label{ch29:def-bm}
A \emph{standard Brownian motion} (or Wiener process) is a random function
$\{W(t): t\ge 0\}$ with these four properties:
\begin{enumerate}
  \item \textbf{Starts at zero:} $W(0) = 0$.
  \item \textbf{Independent increments:} for any times
  $0 \le t_0 < t_1 < \dots < t_k$, the increments
  $W(t_1)-W(t_0),\ W(t_2)-W(t_1),\ \dots$ are mutually independent.
  \item \textbf{Gaussian, stationary increments:} for $s<t$,
  \[
    W(t)-W(s) \;\sim\; \Normal\!\big(0,\; t-s\big).
  \]
  The increment is normal, mean zero, and its variance equals the
  \emph{elapsed time}.
  \item \textbf{Continuous paths:} with probability one, the function
  $t\mapsto W(t)$ is continuous.
\end{enumerate}
\end{definition}

Every one of these falls straight out of the random-walk limit. Independence of
increments is just independence of the coin flips in disjoint time blocks.
Stationarity (the increment law depends only on the gap $t-s$, not on where you
start) is because every coin flip is identical. The Gaussian shape is the Central
Limit Theorem applied to each block of flips. The variance equals elapsed time
because we chose the scaling $\Delta x = 1/\sqrt n$ precisely to make it so. And
continuity is the residue of the steps shrinking to zero.

Notice what properties~1--3 mean for prediction: Brownian motion is a
continuous-time \emph{martingale} (a process modeling a fair game where the best
forecast of future value is its present position). Let $\mathcal{F}_s$ denote the
history of the path up to time $s$. Because the increment $W(t)-W(s)$ is independent
of $\mathcal{F}_s$ and has mean zero, the conditional expectation for any $s < t$ is:
\[
  \E\big[W(t) \mid \mathcal{F}_s\big]
  \;=\; \E\big[W(s) + (W(t)-W(s)) \mid \mathcal{F}_s\big]
  \;=\; W(s) + \E\big[W(t)-W(s)\big]
  \;=\; W(s).
\]
Knowing every twist and turn the path took yesterday gives zero hint whether it
will wander up or down tomorrow: its current state $W(s)$ is its own best forecast.

Brownian motion also carries a second fundamental martingale. Consider
$M(t) = W(t)^2 - t$. Because $\Var(W(t)-W(s)) = t-s$, expanding the square shows
that $\E\big[W(t)^2 - t \mid \mathcal{F}_s\big] = W(s)^2 - s$. The subtracted $t$
neutralizes the deterministic growth of variance, preserving the balance of a
fair game.

Figure~\ref{ch29:fig-martingale} makes this martingale structure visual. Up to
time $s=1$, the path has traced out a specific history to $W(1)=0.6$. Looking
forward, the conditional expectation remains completely flat at $0.6$, while
possible futures fan out symmetrically inside a diffusive trumpet bounded by
$W(s)\pm\sqrt{t-s}$.

\begin{figure}[H]
\centering
\begin{tikzpicture}
\begin{axis}[
  width=12cm, height=6.2cm,
  xlabel={time $t$}, ylabel={position},
  title style={font=\small},
  title={The martingale property: future paths branch around the current state $W(s)$},
  legend pos=north west, legend cell align=left,
  grid=both, grid style={gray!15},
  xmin=0, xmax=2.5, ymin=-0.8, ymax=2.9,
  legend style={font=\footnotesize},
]
\addplot[dashed, thick, gray!70, domain=1:2.5, samples=30, forget plot] {0.6 + sqrt(x-1)};
\addplot[dashed, thick, gray!70, domain=1:2.5, samples=30] {0.6 - sqrt(x-1)};
\addlegendentry{diffusive envelope $W(s)\pm (t-s)^{1/2}$}
\addplot[dashed, very thick, black] coordinates {(1.0,0.6)(2.5,0.6)};
\addlegendentry{forecast $W(s)$}
\addplot[sourceblue, very thick] coordinates {
 (0,0)(0.1,0.2)(0.2,0.1)(0.3,0.35)(0.4,0.25)(0.5,0.45)(0.6,0.3)(0.7,0.5)(0.8,0.4)(0.9,0.55)(1.0,0.6)};
\addlegendentry{path ($s \le 1$)}
\addplot[quotientorange, thick] coordinates {
 (1.0,0.6)(1.2,0.85)(1.4,0.72)(1.6,1.05)(1.8,1.25)(2.0,1.15)(2.2,1.38)(2.5,1.52)};
\addlegendentry{future paths}
\addplot[targetgreen, thick, forget plot] coordinates {
 (1.0,0.6)(1.2,0.52)(1.4,0.68)(1.6,0.55)(1.8,0.64)(2.0,0.56)(2.2,0.65)(2.5,0.62)};
\addplot[densityred, thick, forget plot] coordinates {
 (1.0,0.6)(1.2,0.40)(1.4,0.48)(1.6,0.22)(1.8,0.12)(2.0,-0.15)(2.2,-0.05)(2.5,-0.32)};
\draw[dotted, thick, gray!60] (axis cs:1,-0.8) -- (axis cs:1,2.1);
\addplot[only marks, mark=*, mark size=2.5pt, black] coordinates {(1,0.6)};
\node[anchor=south east, font=\small, inner sep=2pt, fill=white] at (axis cs:0.95,0.68) {$W(s)$};
\end{axis}
\end{tikzpicture}
\caption{The martingale property in picture form. Given the path's history up to
time $s=1$ (solid blue), the expected future position $\E[W(t)\mid \mathcal{F}_s]$
remains flat at the present value $W(s)=0.6$ (black dashed line). Individual
realized trajectories (colored curves) fan out symmetrically, their standard
deviation expanding as $\sqrt{t-s}$ along the diffusive envelope (gray dashed).}
\label{ch29:fig-martingale}
\end{figure}

\begin{focusbox}
The single most important formula in this chapter is
\[
  \boxed{\;\Var\big(W(t)-W(s)\big) = t-s\;}\qquad\text{and in particular}\qquad
  \Var\big(W(t)\big) = t,\quad \SD\big(W(t)\big)=\sqrt t.
\]
Variance grows linearly in time; spread grows like $\sqrt t$. Everything else is
commentary on this fact.
\end{focusbox}

From the definition we can read off the \emph{covariance structure}, which is all
you ever need for a Gaussian process. Take $s\le t$. Write $W(t) = W(s) +
\big(W(t)-W(s)\big)$. The second piece is independent of $W(s)$ and has mean
zero, so
\[
  \Cov\big(W(s),W(t)\big)
  = \Cov\big(W(s),\,W(s)\big) + \Cov\big(W(s),\,W(t)-W(s)\big)
  = \Var(W(s)) + 0 = s.
\]
For general $s,t$ this reads
\[
  \boxed{\;\Cov\big(W(s),W(t)\big) = \min(s,t).\;}
\]

\begin{example}[Working with the increments]
\label{ch29:ex-increments}
Let $W$ be standard Brownian motion. Find $\P\big(W(4) > 3\big)$ and
$\P\big(W(4) - W(1) > 3\big)$.

Since $W(4)\sim\Normal(0,4)$, its standard deviation is $2$, so
\[
  \P(W(4)>3) = \P\!\left(Z > \tfrac{3}{2}\right) = 1-\Phi(1.5)\approx 0.0668,
\]
where $Z$ is standard normal and $\Phi$ its CDF. For the increment,
$W(4)-W(1)\sim\Normal(0,\,4-1)=\Normal(0,3)$, with standard deviation $\sqrt3
\approx 1.732$, so
\[
  \P\big(W(4)-W(1)>3\big) = \P\!\left(Z > \tfrac{3}{\sqrt3}\right)
  = \P(Z>\sqrt3) = 1-\Phi(1.732)\approx 0.0416.
\]
The same threshold of $3$ is less likely to be crossed by a $3$-unit increment
than by the $4$-unit total, because the increment has had less time to spread.
\end{example}

\begin{example}[Covariance in action]
\label{ch29:ex-cov}
Compute $\E\big[W(2)\,W(5)\big]$ and $\Var\big(W(5)-W(2)\big)$.

Means are zero, so $\E[W(2)W(5)] = \Cov(W(2),W(5)) = \min(2,5) = 2$. For the
variance, $W(5)-W(2)\sim\Normal(0,5-2)$, hence $\Var = 3$. Check it the long way:
$\Var(W(5)-W(2)) = \Var(W(5)) - 2\Cov(W(5),W(2)) + \Var(W(2)) = 5 - 2(2) + 2 =
3.$ Same answer.
\end{example}

\begin{recoverybox}
\textbf{Pause and think.} Suppose you track a standard Brownian motion up to time
$s=2$ and observe $W(2) = -1.5$.
\begin{enumerate}[label=(\alph*)]
  \item What is your best forecast for its position at time $t=6$, namely
  $\E[W(6) \mid W(2) = -1.5]$?
  \item What is the conditional standard deviation of $W(6)$ given this observation?
  \item Is the particle more likely to be positive or negative at time $t=6$?
\end{enumerate}
\textit{Check your intuition:}
(a) By the martingale property, $\E[W(6) \mid W(2) = -1.5] = -1.5$. Past moves are
water under the bridge, and future increments have mean zero.
(b) Because increments are independent,
$\Var(W(6) \mid W(2) = -1.5) = \Var(W(6)-W(2)) = 6 - 2 = 4$, so the standard
deviation is $\sqrt{4} = 2$.
(c) Given $W(2) = -1.5$, the position at $t=6$ is distributed as
$\Normal(-1.5, 4)$. Because its mean is negative, it is much more likely to remain
negative: $\P(W(6) > 0 \mid W(2)=-1.5) = \P(Z > 1.5/2) = 1 - \Phi(0.75) \approx 0.2266$,
giving a $77.34\%$ chance of remaining negative.
\end{recoverybox}

\begin{connectionbox}
Brownian motion is a \emph{Gaussian process}: every finite collection
$\big(W(t_1),\dots,W(t_k)\big)$ is a multivariate normal vector. Its mean is the
zero vector and its covariance matrix has entries $\Sigma_{ij}=\min(t_i,t_j)$.
Once you know mean and covariance, a Gaussian object is \emph{completely}
determined -- a fact from the chapter on the multivariate normal and the
covariance matrix in the companion linear-algebra material. So in a real sense
the function $\min(s,t)$ \emph{is} Brownian motion. Sampling a path is just
drawing from a (very high-dimensional) multivariate normal.
\end{connectionbox}

\section{Scaling, Symmetry, and Self-Similarity}
\label{ch29:scaling}

Brownian motion has a family of symmetries that make it feel less like a specific
random scribble and more like a law of nature. Each follows from the covariance
$\min(s,t)$ or directly from the definition.

\begin{proposition}[Invariances]
\label{ch29:prop-invariance}
Let $W$ be standard Brownian motion. Then each of the following is again a
standard Brownian motion:
\begin{enumerate}
  \item \textbf{Scaling (self-similarity):} for any $c>0$,
  $\ \widetilde W(t) = \dfrac{1}{\sqrt c}\,W(ct).$
  \item \textbf{Symmetry:} $\ -W(t).$
  \item \textbf{Time shift (Markov restart):} $\ W(t+s)-W(s)$ for fixed $s\ge0$.
  \item \textbf{Time inversion:} $\ tW(1/t)$ for $t>0$, with value $0$ at $t=0$.
\end{enumerate}
\end{proposition}

\begin{whybox}
Take the scaling property, the most important one. We check it has the right
covariance. With $\widetilde W(t) = c^{-1/2}W(ct)$,
\[
  \Cov\big(\widetilde W(s),\widetilde W(t)\big)
  = \frac1c\,\Cov\big(W(cs),W(ct)\big)
  = \frac1c\,\min(cs,ct)
  = \min(s,t).
\]
Same covariance, still Gaussian, still mean zero, still continuous, still starts
at $0$ -- so it is the same process in law. If you scale
time by a factor $c$ and space by $\sqrt c$, the process is statistically
identical. Brownian motion is self-similar. There is no scale at which it
looks smooth; magnify any segment and you observe the same statistical roughness.
This reflects the $\Delta x = 1/\sqrt n$ scaling from \S\ref{ch29:walk-to-curve}
as an exact symmetry of the limiting process.
\end{whybox}

\begin{aestheticbox}
Self-similarity is why no plotter can ever truly draw Brownian motion. Whatever
resolution you choose, the true path has structure finer than your pixels, with
exactly the same character as what you can see. A drawn Brownian path is always a
\emph{coarsening} -- which is fine, because by self-similarity the coarsening
looks just like the real thing.
\end{aestheticbox}

\section{Continuous Everywhere, Smooth Nowhere}
\label{ch29:nowhere-diff}

Brownian motion possesses an unusual analytic property: its paths are continuous
everywhere, yet differentiable nowhere. The path has no tangent line at any point.

\begin{happeningbox}
\textbf{Why no derivative can exist.} The derivative at $t$, if it existed, would
be the limit of the difference quotient
\[
  \frac{W(t+h)-W(t)}{h}.
\]
But the numerator is $\Normal(0,h)$, so its typical size is $\sqrt h$. Dividing by
$h$ gives a typical size of $\sqrt h / h = 1/\sqrt h$, which \emph{blows up} as
$h\to0$. The increments shrink like $\sqrt h$, but the time gap shrinks faster,
like $h$. The slope is forever losing the race. Over any tiny interval the path
moves a distance proportional to $\sqrt h$, which dwarfs $h$ -- so locally the
graph is far steeper than any line. Zoom in and it does not flatten out into a
tangent the way a smooth curve does; it stays just as rough. That mismatch
between the $\sqrt h$ of space and the $h$ of time is, once again, the diffusive
scaling, and it is exactly what forbids a derivative.
\end{happeningbox}

There is a clean way to quantify the roughness: \emph{quadratic variation}.
Chop $[0,t]$ into $n$ equal pieces of length $\Delta = t/n$ and add up the squared
increments,
\[
  Q_n = \sum_{i=1}^{n}\Big(W(t_i)-W(t_{i-1})\Big)^2 .
\]
Each squared increment has expectation $\Delta$ (since the increment is
$\Normal(0,\Delta)$ and $\E[\text{normal}^2]=\text{variance}$). There are $n$ of
them, so $\E[Q_n] = n\Delta = t$.

Why does $Q_n$ freeze onto the deterministic number $t$ rather than fluctuating
randomly? We can compute its variance directly. The increments
$\Delta W_i = W(t_i) - W(t_{i-1}) \sim \Normal(0, \Delta)$ are mutually independent.
Each can be written as $\Delta W_i = \sqrt{\Delta}\,Z_i$ where $Z_i \sim \Normal(0,1)$.
The squared increment is $(\Delta W_i)^2 = \Delta\,Z_i^2$. For a standard normal,
$\E[Z^2] = 1$ and $\E[Z^4] = 3$, so $\Var(Z^2) = \E[Z^4] - (\E[Z^2])^2 = 3 - 1 = 2$.
Therefore, each squared increment has variance
\[
  \Var\Big((\Delta W_i)^2\Big) = \Var\big(\Delta\,Z_i^2\big) = \Delta^2\,\Var(Z_i^2) = 2\Delta^2.
\]
Because the increments occur over disjoint time intervals, they are independent,
and variances add:
\[
  \Var(Q_n) = \sum_{i=1}^n \Var\Big((\Delta W_i)^2\Big) = n \cdot 2\Delta^2 = 2n\left(\frac{t}{n}\right)^2 = \frac{2t^2}{n}.
\]
As $n\to\infty$, this variance vanishes like $1/n$! By Chebyshev's inequality, the
random sum $Q_n$ concentrates sharply around its mean:
\[
  \boxed{\;Q_n \pto t \quad\text{as } n\to\infty.\;}
\]
This is startling. For an ordinary smooth function $f$, the squared increments are
of order $(\Delta)^2$ and their sum vanishes; for Brownian motion the sum
converges to a \emph{nonzero} limit, the elapsed time itself. This relation is summarized by the formal rule
\[
  (\dd W)^2 = \dd t,
\]
which becomes the foundation of stochastic calculus later in this chapter. By contrast the
ordinary (first-order) variation $\sum |W(t_i)-W(t_{i-1})|$ is \emph{infinite}:
the path wiggles so much that its total arc length over any interval is unbounded.

Figure~\ref{ch29:fig-variation} exposes this deep contrast side by side. As we refine the partition ($n = 8 \to 32 \to 128$), the first-order variation blows up toward infinity because the microscopic step sizes $|\Delta W| \sim 1/\sqrt{n}$ accumulate to $n \cdot (1/\sqrt{n}) = \sqrt{n} \to \infty$. But the quadratic variation freezes steadily into the deterministic diagonal line $Q(t) = t$. Although each individual step is random, the fluctuations of $(\Delta W_i)^2$ around its mean $\Delta$ average out across the fine partition, leaving the deterministic limit $t$.

\begin{figure}[H]
\centering
\begin{tikzpicture}
\begin{axis}[
  name=leftplot,
  width=6.3cm, height=5.5cm,
  xlabel={time $t$}, ylabel={quadratic variation $Q_n(t)$},
  title style={font=\small},
  title={Quadratic variation: $Q_n(t) \to t$},
  legend pos=north west, legend cell align=left,
  grid=both, grid style={gray!15},
  xmin=0, xmax=1, ymin=0, ymax=1.3,
  legend style={font=\footnotesize},
]
\addplot[dashed, thick, black] coordinates {(0,0)(1,1)};
\addlegendentry{limit $y=t$}
\addplot[thick, sourceblue] coordinates {
 (0,0)(0.125,0.08)(0.25,0.31)(0.375,0.39)(0.5,0.48)(0.625,0.72)(0.75,0.79)(0.875,0.91)(1.0,1.06)};
\addlegendentry{$n=8$}
\addplot[thick, quotientorange] coordinates {
 (0,0)(0.1,0.09)(0.2,0.22)(0.3,0.28)(0.4,0.39)(0.5,0.51)(0.6,0.58)(0.7,0.72)(0.8,0.79)(0.9,0.92)(1.0,1.01)};
\addlegendentry{$n=32$}
\addplot[very thick, targetgreen] coordinates {
 (0,0)(0.1,0.10)(0.2,0.20)(0.3,0.30)(0.4,0.40)(0.5,0.50)(0.6,0.60)(0.7,0.70)(0.8,0.80)(0.9,0.90)(1.0,1.00)};
\addlegendentry{$n=128$}
\end{axis}

\begin{axis}[
  at={(leftplot.south east)}, xshift=2.6cm,
  width=6.3cm, height=5.5cm,
  xlabel={time $t$}, ylabel={total variation $V_1(t)$},
  ylabel style={yshift=2pt},
  title style={font=\small},
  title={Total variation: $V_1(t) \to \infty$},
  legend pos=north west, legend cell align=left,
  grid=both, grid style={gray!15},
  xmin=0, xmax=1, ymin=0, ymax=12,
  legend style={font=\footnotesize},
]
\addplot[thick, sourceblue] coordinates {
 (0,0)(0.125,0.28)(0.25,0.68)(0.375,0.96)(0.5,1.26)(0.625,1.75)(0.75,2.01)(0.875,2.36)(1.0,2.65)};
\addlegendentry{$n=8$}
\addplot[thick, quotientorange] coordinates {
 (0,0)(0.1,0.52)(0.2,1.05)(0.3,1.55)(0.4,2.08)(0.5,2.60)(0.6,3.12)(0.7,3.65)(0.8,4.18)(0.9,4.71)(1.0,5.20)};
\addlegendentry{$n=32$}
\addplot[very thick, densityred] coordinates {
 (0,0)(0.1,1.02)(0.2,2.05)(0.3,3.08)(0.4,4.11)(0.5,5.13)(0.6,6.16)(0.7,7.19)(0.8,8.21)(0.9,9.24)(1.0,10.25)};
\addlegendentry{$n=128$}
\end{axis}
\end{tikzpicture}
\caption{The dual nature of Brownian roughness. Left: the quadratic variation
$Q_n(t) = \sum_{t_i \le t} (\Delta W_i)^2$ collapses cleanly onto the deterministic line
$y=t$ as $n\to\infty$. Right: the total arc length $V_1(t) = \sum_{t_i \le t} |\Delta W_i|$
diverges like $\sqrt{n}$ into an infinitely jagged fractal path.}
\label{ch29:fig-variation}
\end{figure}

\begin{trapbox}
Do not picture the velocity of a pollen grain as a Brownian path. A real grain has
a real velocity at every instant -- Brownian motion is an idealization that
ignores the grain's inertia, valid only on time scales longer than the gap
between collisions. The mathematical object $W(t)$ has \emph{no velocity}, which
is fine because it models \emph{position} on time scales coarse enough that
individual molecular velocities have averaged out. Confusing the model with the physical grain leads
to paradoxes like ``infinite kinetic energy'' that are artifacts of the
idealization.
\end{trapbox}

\begin{paradoxbox}
\textbf{A continuous function with no derivative anywhere?} Before the late
1800s, it was widely believed that a continuous function must be
differentiable except perhaps at isolated points. In 1872, Karl
Weierstrass constructed an explicit continuous-but-nowhere-differentiable function,
which Charles Hermite and many contemporaries dismissed as a pathological curiosity. Brownian motion demonstrates that this behavior arises naturally in stochastic systems:
a standard Brownian path is nowhere differentiable with probability one. In the space of continuous functions,
paths with bounded variation and continuous derivatives are the exception rather than the rule.
\end{paradoxbox}

\section{The Heat Equation: Randomness as Diffusion}
\label{ch29:heat}

Now comes a striking link: this random, jagged process is governed by a smooth
deterministic partial differential equation -- the \emph{heat equation}. The same
equation describes heat spreading through a
metal bar, ink diffusing in still water, and the probability density of a
Brownian particle.

Let $p(x,t)$ be the probability density of $W(t)$. From the definition,
$W(t)\sim\Normal(0,t)$, so
\[
  p(x,t) = \frac{1}{\sqrt{2\pi t}}\,\exp\!\left(-\frac{x^2}{2t}\right).
\]
This is the Gaussian, but with its variance \emph{growing in time}. At $t$ near
$0$ it is a tall thin spike at the origin (the particle is almost surely near
$0$); as $t$ grows the bump flattens and spreads, exactly like a drop of dye
diffusing outward.

\begin{figure}[H]
\centering
\begin{tikzpicture}
\begin{axis}[
  width=12cm, height=6cm,
  xlabel={position $x$}, ylabel={density $p(x,t)$},
  title style={font=\small},
  title={The density of $W(t)$ spreads and flattens: a diffusing Gaussian},
  domain=-5:5, samples=200, legend pos=north east, legend cell align=left,
  grid=both, grid style={gray!15}, ymin=0,
]
\addplot[densityred, very thick] {1/sqrt(2*pi*0.5)*exp(-x^2/(2*0.5))};
\addlegendentry{$t=0.5$}
\addplot[quotientorange, thick] {1/sqrt(2*pi*1.5)*exp(-x^2/(2*1.5))};
\addlegendentry{$t=1.5$}
\addplot[sourceblue, thick] {1/sqrt(2*pi*4)*exp(-x^2/(2*4))};
\addlegendentry{$t=4$}
\end{axis}
\end{tikzpicture}
\caption{Snapshots of the Brownian density at three times. The area under each
curve stays $1$ (it is always a probability density); the bump just spreads. This
is the same evolution heat undergoes in a conducting bar.}
\label{ch29:fig-heat}
\end{figure}

Now differentiate. A short calculation (done in the box below) shows that this
density satisfies
\[
  \boxed{\;\frac{\partial p}{\partial t} = \frac{1}{2}\frac{\partial^2 p}{\partial x^2}.\;}
\]
This is the \emph{heat equation} (with diffusion constant $\tfrac12$). The
left side is the rate at which the density changes in time; the right side is its
curvature in space. Where the density is concave down (a peak), the right side is
negative and the peak sinks; where it is concave up (a valley between two bumps),
probability fills in. Probability flows from high to low, smoothing itself out --
the exact behavior of temperature.

\begin{whybox}
Verify $\partial_t p = \tfrac12 \partial_{xx} p$ directly for
$p = (2\pi t)^{-1/2}e^{-x^2/2t}$. Take logs to make differentiation easy:
$\ln p = -\tfrac12\ln(2\pi t) - \tfrac{x^2}{2t}.$ Then
\[
  \frac{\partial_t p}{p} = -\frac{1}{2t} + \frac{x^2}{2t^2},
  \qquad
  \frac{\partial_x p}{p} = -\frac{x}{t}.
\]
For the second space derivative,
$\partial_{xx} p = \partial_x\!\big(p\cdot(-x/t)\big)
= p\big(-x/t\big)\big(-x/t\big) + p\big(-1/t\big)
= p\big(x^2/t^2 - 1/t\big).$ Therefore
\[
  \tfrac12\,\partial_{xx} p = p\left(\frac{x^2}{2t^2}-\frac{1}{2t}\right)
  = \partial_t p,
\]
matching the time derivative line for line. The Gaussian is not just \emph{a}
solution of the heat equation; with a point-mass initial condition it is \emph{the}
solution, the so-called heat kernel.
\end{whybox}

\begin{happeningbox}
\textbf{Why do randomness and heat obey the same equation?} Because both are made
of independent little pushes. Heat is energy random-walking from molecule to
molecule; a Brownian particle is position random-walking under molecular kicks.
The heat equation is just the continuum bookkeeping of ``each point passes a bit
of its stuff to its neighbors on each side.'' Whether the stuff is energy or
probability, the local rule is identical, so the macroscopic law is identical. The
factor of $\tfrac12$ is the fingerprint of the variance-equals-time normalization
in Definition~\ref{ch29:def-bm}; a different step variance would rescale it.
\end{happeningbox}

\begin{connectionbox}
This is the gateway to the \emph{Feynman--Kac formula} and to physics. Solving a
diffusion PDE can be turned into averaging a function over Brownian paths, and
vice versa. Expectation over random paths $\leftrightarrow$ solution of a
deterministic PDE. The same duality, rotated into imaginary time, connects
Brownian motion to the Schr\"odinger equation of quantum mechanics. The integral
sign and the expectation sign are, here, two views of one object -- a theme the
companion calculus book develops as ``expectation is an integral.''
\end{connectionbox}

\section{First Passages, the Reflection Principle, and the Maximum}
\label{ch29:reflection}

A natural question: starting from $0$, when does Brownian motion first reach a
level $a>0$? Call that random time $\tau_a = \inf\{t: W(t)=a\}$. By continuity the
path cannot jump over $a$; it must pass through it. The \emph{reflection principle}
provides an exact solution using path symmetry.

\begin{whybox}
\textbf{The reflection principle.} Fix a level $a>0$ and a time $t$. We want
$\P\big(\max_{0\le s\le t} W(s) \ge a\big)$, the chance the path ever reaches $a$
by time $t$. Split the event $\{\max_{0\le s\le t} W(s)\ge a\}$: any path ending above $a$ has certainly
hit $a$ on the way (continuity again). Now take any path that hits $a$ at some
first time $\tau\le t$ but ends \emph{below} $a$, and reflect its portion after
$\tau$ across the line $a$. By the symmetry of Brownian motion (after $\tau$ it
restarts fresh, and up is as likely as down), the reflected path is exactly as
probable as the original, and it ends \emph{above} $a$. This pairs up each
``hit-then-below'' path with an equally likely ``hit-then-above'' path. Hence
\[
  \P\Big(\max_{0\le s\le t} W(s)\ge a\Big) = 2\,\P\big(W(t)>a\big).
\]
The maximum reaching $a$ is exactly twice as likely as the endpoint exceeding $a$.
\end{whybox}

\begin{figure}[H]
\centering
\begin{tikzpicture}
\begin{axis}[
  width=12cm, height=6cm,
  xlabel={time $t$}, ylabel={position},
  title style={font=\small},
  title={The reflection principle: each path below $a$ pairs with a mirror image above},
  legend pos=south west, legend cell align=left,
  grid=both, grid style={gray!15},
  ymin=-0.7, ymax=2.1, xmin=0, xmax=1,
]
% barrier level a=1.2
\addplot[black, dashed, thick, domain=0:1, samples=2] {1.2};
\addlegendentry{barrier $a$}
% original path: hits a at tau~0.5, then ends below a
\addplot[sourceblue, thick] coordinates {
 (0,0)(0.08,0.25)(0.16,0.12)(0.24,0.45)(0.32,0.7)(0.4,0.95)(0.5,1.2)
 (0.58,1.0)(0.66,1.15)(0.74,0.85)(0.82,0.95)(0.9,0.7)(1,0.6)};
\addlegendentry{original (hits $a$, ends below)}
% reflected path: same up to tau, then mirrored across a after 0.5
\addplot[densityred, thick, dashed] coordinates {
 (0.5,1.2)(0.58,1.4)(0.66,1.25)(0.74,1.55)(0.82,1.45)(0.9,1.7)(1,1.8)};
\addlegendentry{reflected (ends above)}
% mark hitting time
\addplot[only marks, mark=*, mark size=2pt, black] coordinates {(0.5,1.2)};
\node[anchor=south east] at (axis cs:0.48,1.22) {\small $\tau$};
\end{axis}
\end{tikzpicture}
\caption{After the first hitting time $\tau$, reflect the path across the level $a$.
Because Brownian motion restarts fresh and symmetrically at $\tau$, the reflected
post-$\tau$ piece (red, ending above $a$) is exactly as probable as the original
(blue, ending below $a$). This pairing is why
$\P(\max_{[0,t]}W\ge a)=2\,\P(W(t)>a)$.}
\label{ch29:fig-reflection}
\end{figure}

Figure~\ref{ch29:fig-reflection} shows the pairing in action. Since
$W(t)\sim\Normal(0,t)$, this gives the clean formula
\[
  \P\Big(\max_{0\le s\le t}W(s)\ge a\Big)
  = 2\,\P\big(W(t)>a\big)
  = 2\Big(1-\Phi\big(a/\sqrt t\big)\Big).
\]

From this maximum probability, we immediately obtain the exact distribution and density of the first hitting time $\tau_a = \inf\{t > 0 : W(t) = a\}$. Because Brownian paths are continuous, reaching or exceeding level $a$ by time $t$ is logically equivalent to hitting $a$ at or before time $t$:
\[
  \{\tau_a \le t\} \;=\; \Big\{\max_{0\le s\le t} W(s) \ge a\Big\}.
\]
Therefore, the cumulative distribution function of $\tau_a$ is
\[
  F_{\tau_a}(t) = \P(\tau_a \le t) = 2\Big(1 - \Phi\big(a/\sqrt{t}\big)\Big) = \frac{2}{\sqrt{2\pi}} \int_{a/\sqrt{t}}^\infty e^{-u^2/2}\,\dd u.
\]
To find the probability density function $f_{\tau_a}(t)$, differentiate with respect to $t$ via the Leibniz rule:
\[
  f_{\tau_a}(t) = \frac{\dd}{\dd t} F_{\tau_a}(t)
  = -\frac{2}{\sqrt{2\pi}} \exp\!\left(-\frac{a^2}{2t}\right) \cdot \frac{\dd}{\dd t}\left(\frac{a}{\sqrt{t}}\right)
  = -\frac{2}{\sqrt{2\pi}} \exp\!\left(-\frac{a^2}{2t}\right) \left(-\frac{a}{2 t^{3/2}}\right).
\]
Simplifying the signs and factors of $2$ yields the closed-form density
\[
  \boxed{\;f_{\tau_a}(t) = \frac{a}{\sqrt{2\pi t^3}}\,\exp\!\left(-\frac{a^2}{2t}\right), \qquad t > 0.\;}
\]
This is the celebrated \emph{L\'evy distribution} (an inverse Gaussian). Notice its heavy power-law tail: for large $t$, $\exp(-a^2/(2t)) \to 1$, so $f_{\tau_a}(t) \sim \frac{a}{\sqrt{2\pi}}\,t^{-3/2}$. As a consequence, although the particle is guaranteed to hit level $a$ eventually ($\lim_{t\to\infty} F_{\tau_a}(t) = 2(1 - \Phi(0)) = 1$), its expected hitting time is infinite:
\[
  \E[\tau_a] = \int_0^\infty t\,f_{\tau_a}(t)\,\dd t \;\approx\; \frac{a}{\sqrt{2\pi}} \int^\infty t^{-1/2}\,\dd t = \infty.
\]
The particle wanders into endlessly deep excursion loops away from $a$, so its average waiting time blows up even though it never escapes forever.

\begin{example}[Hitting a barrier]
\label{ch29:ex-barrier}
A particle does standard Brownian motion. What is the probability it reaches level
$a=2$ at some point during $[0,4]$?

Here $a/\sqrt t = 2/\sqrt4 = 1$, so
\[
  \P\Big(\max_{0\le s\le4}W(s)\ge 2\Big) = 2\big(1-\Phi(1)\big)
  = 2(0.1587) = 0.3174.
\]
Even though the \emph{endpoint} $W(4)$ exceeds $2$ only with probability
$1-\Phi(1)=0.1587$, the path \emph{touches} $2$ at some moment with twice that
chance, because it can hit the barrier and wander back down.
\end{example}

\begin{aestheticbox}
The reflection principle is a rare case where a hard analytic question -- the
distribution of a maximum over a whole continuum of times -- collapses to a
one-line symmetry argument. It works because of three facts stacked together:
continuity (no jumping over $a$), the strong Markov property (a fresh restart at
the hitting time), and up-down symmetry. Lose any one and the trick breaks.
\end{aestheticbox}

\section{Variations: Drift and Geometric Brownian Motion}
\label{ch29:drift}

Pure Brownian motion has no preferred direction. Real systems often do: a stock
tends to drift upward, a diffusing particle in a current is carried along. We add
a deterministic trend and a volatility scale:
\[
  X(t) = \mu t + \sigma W(t).
\]
Here $\mu$ is the \emph{drift} (mean velocity) and $\sigma$ the \emph{volatility}.
Then $X(t)\sim\Normal(\mu t,\ \sigma^2 t)$: the mean marches linearly while the
spread still grows like $\sqrt t$. For small $t$ the random $\sqrt t$ spread
dominates the linear $\mu t$ trend; for large $t$ the trend wins. That crossover
is why short-term price moves look like pure noise but long-term ones show the
drift.

For prices we want something that cannot go negative and whose \emph{percentage}
changes, not absolute changes, behave nicely. The standard model is
\emph{geometric Brownian motion},
\[
  S(t) = S(0)\,\exp\!\Big(\big(\mu - \tfrac12\sigma^2\big)t + \sigma W(t)\Big).
\]
The log-price is Brownian motion with drift, so the price itself is the
exponential of a normal -- a \emph{lognormal} random variable. The extra
$-\tfrac12\sigma^2$ drift term is genuine: it comes from Jensen's inequality and
the quadratic-variation rule $(\dd W)^2=\dd t$, as we now verify.

\begin{figure}[H]
\centering
\begin{tikzpicture}
\begin{axis}[
  width=12cm, height=5.6cm,
  xlabel={time $t$}, ylabel={value},
  title style={font=\small},
  title={Brownian motion with drift (blue) and a geometric path (orange)},
  legend pos=north west, legend cell align=left, grid=both, grid style={gray!15},
  ymin=-0.1, ymax=2.2,
]
\addplot[sourceblue, thick] coordinates {
 (0,0)(0.1,0.18)(0.2,0.12)(0.3,0.34)(0.4,0.30)(0.5,0.48)(0.6,0.42)
 (0.7,0.66)(0.8,0.60)(0.9,0.84)(1,0.95)};
\addlegendentry{$X(t)=\mu t+\sigma W(t)$}
\addplot[quotientorange, thick] coordinates {
 (0,1)(0.1,1.05)(0.2,1.02)(0.3,1.14)(0.4,1.10)(0.5,1.24)(0.6,1.18)
 (0.7,1.36)(0.8,1.30)(0.9,1.52)(1,1.60)};
\addlegendentry{$S(t)$ (geometric)}
\addplot[densityred, dashed, thick, domain=0:1, samples=2] {0.9*x};
\addlegendentry{mean trend $\mu t$}
\end{axis}
\end{tikzpicture}
\caption{Adding a drift tilts the path; the geometric version exponentiates it so
the value stays positive and compounds. Both still carry the unmistakable
Brownian jitter.}
\label{ch29:fig-drift}
\end{figure}

\section{A First Look at Stochastic Calculus}
\label{ch29:ito}

Because Brownian paths are nowhere differentiable, the classical derivative
$\dd W/\dd t$ does not exist. Instead, calculus with Brownian motion is developed in
\emph{integral} form through \emph{It\^o calculus}, built around the key relation
$(\dd W)^2 = \dd t$.

The central object is the It\^o integral $\int_0^t H(s)\,\dd W(s)$, which you
should read as the limiting payoff of a strategy that holds $H(s)$ units of a
Brownian asset and collects the increments $\dd W$. The key feature is that the integrand is
evaluated at the \emph{left} endpoint of each tiny interval -- you must decide
your holding \emph{before} the next random increment, not after. That ``no
peeking into the future'' rule is what keeps the theory well behaved and
tractable, and it is what links It\^o integrals to martingales (the previous
chapter).

The headline result is \textbf{It\^o's lemma}, the chain rule of this world. For
an ordinary smooth function $f$ and ordinary calculus, $\dd f(W) = f'(W)\,\dd W$.
But Brownian motion forces you to keep one more term, because $(\dd W)^2$ does not
vanish -- it equals $\dd t$:
\[
  \boxed{\;\dd f\big(W(t)\big) = f'\big(W(t)\big)\,\dd W(t)
  + \tfrac12 f''\big(W(t)\big)\,\dd t.\;}
\]
That extra $\tfrac12 f''\,\dd t$ is the entire difference between ordinary and
stochastic calculus, and it is the quadratic variation of
\S\ref{ch29:nowhere-diff} showing up to collect its due.

\begin{whybox}
Where does the extra term come from? Taylor-expand $f$ over a tiny step:
\[
  f(W+\dd W) - f(W)
  = f'(W)\,\dd W + \tfrac12 f''(W)\,(\dd W)^2 + \cdots.
\]
For a smooth path $(\dd W)^2$ would be negligible compared to $\dd W$, and you
would drop it -- that is ordinary calculus. But for Brownian motion $(\dd W)^2$ is
\emph{not} negligible: summed up it converges to $\dd t$, not to zero. So the
second-order term survives into the answer while the cubic and higher terms still
vanish. It\^o's lemma is just Taylor's theorem with the rule $(\dd W)^2=\dd t$
substituted in. The nowhere-differentiability we feared turns out to be exactly
what we exploit.
\end{whybox}

\begin{example}[Why the $-\tfrac12\sigma^2$ in geometric Brownian motion]
\label{ch29:ex-gbm}
Apply It\^o's lemma to $f(w)=e^{\sigma w}$, so $f'=\sigma e^{\sigma w}$,
$f''=\sigma^2 e^{\sigma w}$. With $Y(t)=e^{\sigma W(t)}$,
\[
  \dd Y = \sigma Y\,\dd W + \tfrac12\sigma^2 Y\,\dd t.
\]
The exponential of Brownian motion drifts \emph{upward} at rate
$\tfrac12\sigma^2$ all by itself, purely from the convexity of $e^x$ feeding on
the wiggle. To model a stock whose true growth rate is $\mu$, you must
\emph{subtract} that drift, which is why
$S(t)=S_0\exp\big((\mu-\tfrac12\sigma^2)t+\sigma W(t)\big)$ carries the correction
term: the $-\tfrac12\sigma^2$ is precisely the second-order adjustment from It\^o's lemma.
\end{example}

\begin{applicationbox}
\textbf{Black--Scholes and the price of an option.} In 1973 Fischer Black, Myron
Scholes, and Robert Merton modeled a stock as geometric Brownian motion and asked:
what is the fair price today of a \emph{call option}, the right to buy the stock
at a fixed price $K$ at a future date $T$? Their insight was to construct a dynamic
hedging portfolio. At each instant you can hold a combination of the stock and cash that
\emph{exactly cancels} the random $\dd W$ term in the option's value.
With the randomness hedged away, the portfolio must (to
preclude arbitrage) earn the risk-free interest rate $r$. Setting those two facts
equal and applying It\^o's lemma transforms the problem into a partial differential equation -- the
Black--Scholes equation -- which can be transformed into the heat equation of \S\ref{ch29:heat}.
Solving it gives the closed-form price
\[
  C = S_0\,\Phi(d_1) - K e^{-rT}\,\Phi(d_2),
  \qquad
  d_{1,2} = \frac{\ln(S_0/K) + (r \pm \tfrac12\sigma^2)T}{\sigma\sqrt T}.
\]
There is that $\tfrac12\sigma^2$ and that $\Phi$ -- the Gaussian CDF -- again. A
nineteenth-century botanist's jittering pollen grain became, a century and a half
later, the foundation of a multitrillion-dollar derivatives industry and a Nobel
Prize in economics. That is the reach of one random curve.
\end{applicationbox}

\begin{trapbox}
Geometric Brownian motion is a \emph{model}, and a famously imperfect one. Real
returns have fatter tails than the lognormal allows, volatility $\sigma$ is not
constant, and prices occasionally jump (Brownian paths never do). The 1987 crash
and 2008 crisis both featured moves that Black--Scholes assigns probabilities so
small they ``should not happen in the lifetime of the universe.'' Use the model
for intuition and as a baseline; respect its assumptions; never mistake the
elegant curve for the messy market.
\end{trapbox}

\section{Simulating Brownian Motion}
\label{ch29:simulate}

The cleanest way to \emph{see} Brownian motion is to generate one. Definition
\ref{ch29:def-bm} hands you the recipe directly: increments over disjoint
intervals are independent normals with variance equal to the time step. So pick a
step $\Delta t$, draw independent $Z_i\sim\Normal(0,1)$, and accumulate.

\begin{labbox}
Build a Brownian path on $[0,1]$ with $N=1000$ steps, so $\Delta t = 1/N$.
\begin{enumerate}
  \item Set $W_0 = 0$.
  \item For $i=1,\dots,N$, draw $Z_i\sim\Normal(0,1)$ and set
  \[
    W_i = W_{i-1} + \sqrt{\Delta t}\; Z_i.
  \]
  The factor $\sqrt{\Delta t}$ is the diffusive scaling again: variance per step
  equals $\Delta t$, so the increment's standard deviation is $\sqrt{\Delta t}$.
  \item Plot $(i\,\Delta t,\ W_i)$.
\end{enumerate}
Now experiment. Halve $\Delta t$ (double $N$) and overlay: the new path looks just
as rough -- self-similarity in your own plot. Replace the update with
$W_i = W_{i-1} + \mu\,\Delta t + \sigma\sqrt{\Delta t}\,Z_i$ to add drift and
volatility. Exponentiate to get a stock-price path. Run it a thousand times and
histogram $W_1$: you will recover the standard normal, the Central Limit Theorem
closing the loop we opened in \S\ref{ch29:walk-to-curve}.
\end{labbox}

\begin{remark}
That update rule, $W_i = W_{i-1} + \sqrt{\Delta t}\,Z_i$, is the
\emph{Euler--Maruyama} scheme, the standard numerical method for stochastic simulation. Note the
$\sqrt{\Delta t}$ where an ordinary differential-equation solver would use
$\Delta t$. Omitting the square root simulates a different process and violates the diffusive scaling.
\end{remark}

\section{What to Carry Away}
\label{ch29:summary}

\begin{focusbox}
\begin{itemize}
  \item Brownian motion is the scaling limit of random walks under the
  \emph{diffusive scaling}: time $\sim 1/n$, space $\sim 1/\sqrt n$. Space is the
  square root of time.
  \item It is defined by four properties: starts at $0$, independent increments,
  Gaussian stationary increments with $W(t)-W(s)\sim\Normal(0,t-s)$, and
  continuous paths. Equivalently it is the Gaussian process with covariance
  $\min(s,t)$.
  \item Its paths are continuous everywhere and differentiable nowhere; their
  quadratic variation over $[0,t]$ equals $t$, encoded as $(\dd W)^2 = \dd t$.
  \item The density of $W(t)$ is a spreading Gaussian that solves the heat
  equation $\partial_t p = \tfrac12\partial_{xx}p$ -- randomness and diffusion are
  the same mathematics.
  \item The reflection principle gives maxima and first-passage probabilities for
  free: $\P(\max_{[0,t]}W\ge a) = 2\,\P(W(t)>a)$.
  \item It\^o's lemma, $\dd f(W) = f'(W)\dd W + \tfrac12 f''(W)\dd t$, is the chain
  rule of a world where $(\dd W)^2=\dd t$, and it powers geometric Brownian motion
  and the Black--Scholes formula.
\end{itemize}
\end{focusbox}

The thread running through every section is a single scaling law -- space grows
like the square root of time -- arising from the cancellation of independent random
steps. From this scaling flow the Gaussian increments, the roughness of the paths,
the lack of differentiability, the connection to the heat equation, and the second-order term in stochastic
calculus. Brown documented the physical phenomenon under a microscope; Einstein formulated its physical theory;
Wiener established its mathematical foundation; and Black, Scholes, and Merton applied it to financial economics.

So far, we took distributions as given: we assumed a coin has bias $p=1/2$, a
queue has arrival rate $\lambda$, or Brownian motion has drift $\mu$, and deduced
how the system behaves. In practice, nobody hands us the true parameters. We
only observe data -- coin flips, arrival times, noisy prices -- and must infer the
mechanism that produced them. Chapter~30 turns probability around to begin
statistics, using the method of moments and maximum likelihood to estimate
unknown parameters from observed data.

\section{Exercises}
\label{ch29:exercises}

\begin{exercise}[Increment arithmetic]
Let $W$ be standard Brownian motion. Compute: (a) $\Var(W(7)-W(3))$;
(b) $\Cov(W(3),W(7))$; (c) $\P(W(9)>6)$ in terms of $\Phi$; and (d) $\E[W(2)^2]$.
\end{exercise}

\begin{exercise}[Correlation between times]
For $s<t$, show that
$\Corr\big(W(s),W(t)\big) = \sqrt{s/t}$. Interpret: why does the correlation
between the position now and the position far in the future decay, and at what
rate?
\end{exercise}

\begin{exercise}[Scaling self-check]
Let $\widetilde W(t) = \tfrac{1}{2}W(4t)$. Verify directly that
$\widetilde W(t)\sim\Normal(0,t)$ and that
$\Cov(\widetilde W(s),\widetilde W(t)) = \min(s,t)$, confirming
$\widetilde W$ is again a standard Brownian motion.
\end{exercise}

\begin{exercise}[Brownian bridge]
Define $B(t) = W(t) - t\,W(1)$ for $0\le t\le1$ (Brownian motion ``tied down'' to
$0$ at both ends). Show $B(0)=B(1)=0$, that $\E[B(t)]=0$, and that
$\Cov\big(B(s),B(t)\big) = \min(s,t) - st$ for $s,t\in[0,1]$.
\end{exercise}

\begin{exercise}[Reflection principle]
A standard Brownian motion runs on $[0,9]$. (a) Find the probability it reaches
level $a=3$ at some time during $[0,9]$. (b) Find the probability it reaches $3$
but ends below $3$, i.e.\ $\P(\max_{[0,9]}W\ge3,\ W(9)<3)$.
(Hint: use the reflection pairing.)
\end{exercise}

\begin{exercise}[Drifted Brownian motion]
Let $X(t)=\mu t + \sigma W(t)$ with $\mu=1$, $\sigma=2$. (a) Give the distribution
of $X(4)$. (b) Compute $\P(X(4)>10)$ in terms of $\Phi$. (c) For what time $t$
does the mean of $X(t)$ first exceed its standard deviation?
\end{exercise}

\begin{exercise}[Heat equation, by hand]
Verify that $p(x,t)=(4\pi t)^{-1/2}\exp(-x^2/4t)$ satisfies
$\partial_t p = \partial_{xx}p$ (diffusion constant $1$ rather than $\tfrac12$).
Which variance growth rate does this correspond to, and how does the diffusion
constant relate to the rate at which $\Var$ grows in time?
\end{exercise}

\begin{exercise}[It\^o's lemma practice]
Use It\^o's lemma to find the stochastic differential of (a) $f(W)=W^2$ and
(b) $f(W)=W^3$. For part (a), take expectations of both sides and integrate to
re-derive $\E[W(t)^2]=t$. (Hint: the It\^o-integral term has expectation zero.)
\end{exercise}

\begin{exercise}[Geometric Brownian motion]
Let $S(t)=S_0\exp\big((\mu-\tfrac12\sigma^2)t+\sigma W(t)\big)$. (a) Show that
$\ln S(t)\sim\Normal\big(\ln S_0 + (\mu-\tfrac12\sigma^2)t,\ \sigma^2 t\big)$.
(b) Using $\E[e^{Z}] = e^{\Var(Z)/2}$ for mean-zero normal $Z$, show that
$\E[S(t)] = S_0 e^{\mu t}$ -- so the $-\tfrac12\sigma^2$ correction makes the
\emph{expected} price grow at exactly rate $\mu$.
\end{exercise}

\begin{exercise}[Conceptual: why no derivative]
Explain in your own words, using the $\sqrt h$ versus $h$ scaling of increments,
why $\big(W(t+h)-W(t)\big)/h$ has no finite limit as $h\to0$. Then explain why the
quadratic variation $\sum (\Delta W)^2$ converges to a nonzero limit while it
would vanish for any differentiable function.
\end{exercise}

\begin{exercise}[Simulation, proof by experiment]
Describe an Euler--Maruyama scheme to simulate $W$ on $[0,T]$ with $N$ steps.
Explain why the per-step increment must use $\sqrt{\Delta t}$ and not $\Delta t$,
and predict what histogram you would get from $W(T)$ over many independent runs.
What does the histogram demonstrate about the connection to the Central Limit
Theorem?
\end{exercise}
